\documentclass[../../main.tex]{subfiles}
\begin{document}

{\bf[解]}

円$C$を含む平面を$xy$平面とし，$O$を原点，$A(0,a,0)$（$a>0$）とする．直線$OA$は$y$軸方向だから，$l$は$OA$と直交し$xy$平面と$45^\circ$をなす直線として，方向ベクトル
$\vec{v}=
\begin{pmatrix}1\\0\\-1
\end{pmatrix}$
がとれる．

\begin{figure}[htb]
  \centering
  \tdplotsetmaincoords{75}{115}
  \begin{tikzpicture}[scale=1.5, >=stealth, every node/.style={font=\small}]
    \begin{scope}[tdplot_main_coords]
      % 円 C（半径1，xy平面内）
      \draw[dashed, gray!55!black] plot[domain=0:360, samples=80, variable=\t]
      ({cos(\t)}, {sin(\t)}, 0);

      % 座標軸
      \draw[->] (-1.8,0,0) -- (1.8,0,0) node[below left=1pt] {$x$};
      \draw[->] (0,-0.3,0) -- (0,2.3,0) node[right] {$y$};
      \draw[->] (0,0,-0.3) -- (0,0,2.3) node[above] {$z$};
      \node[below right=1pt] at (0,0,0) {$O$};

      % 点 A(0,a,0) と直線 l（方向ベクトル (1,0,-1)）
      \coordinate (A) at (0,1.8,0);
      \draw[thick] ($(A)+(-1.3,0,1.3)$) -- ($(A)+(1.3,0,-1.3)$) node[above right=1pt] {$l$};
      \fill (A) circle (1.2pt) node[below right=1pt] {$A$};

      % x軸に並行
      \draw[dashed] ($(A)+(-1.3,0,0)$) -- ($(A)+(1.3,0,0)$);

      % l と xy平面のなす角 45°（x軸方向から l 方向への弧。両者とも A を含む xz 平面内にあるので
      % A+r(cos(phi),0,sin(phi)) が phi=0 で x軸方向，phi=-45度で l の方向と一致する）
      \draw[gray!40!black] plot[domain=0:-45, samples=20, variable=\phi]
      ({0.5*cos(\phi)}, {1.8}, {0.5*sin(\phi)});
      \node[font=\scriptsize] at ($(A)+(0.62,0,-0.22)$) {$45^\circ$};
    \end{scope}
  \end{tikzpicture}
  \caption{円 $C$，点 $A$ を通り直線 $OA$ と直交し $xy$ 平面と $45^\circ$ をなす直線 $l$}
  \label{fig:overview_3d}
\end{figure}

$l$上の点$X$をパラメータ$t\in\mathbb{R}$を用いて
\begin{align}
  \overrightarrow{OX}
  =
  \begin{pmatrix}0\\a\\0
  \end{pmatrix}+t
  \begin{pmatrix}1\\0\\-1
  \end{pmatrix}
  =
  \begin{pmatrix}t\\a\\-t
  \end{pmatrix}  \label{eq:1}
\end{align}
とする．対称性から以下$t\ge0$として考える．
また実数$\theta\, (0\le\theta<2\pi)$を用いて$C$上の点$P$は$P(\cos\theta,\sin\theta,0)$と表せる．

このとき$|PX|^2$は
\begin{align}
  |PX|^2&=(t-\cos\theta)^2+(a-\sin\theta)^2+t^2 \\
  &=(2t^2+a^2+1)-2(t\cos\theta+a\sin\theta) \label{eq:2}
\end{align}
である．まず$t,a$を固定して$\theta$を動かした時の最小値を考える．
$t\ge0,\,a>0$だから，コーシー・シュワルツの不等式より
\begin{align}
  t\cos\theta+a\sin\theta\le\sqrt{t^2+a^2}\sqrt{\cos^2\theta+\sin^2\theta}=\sqrt{t^2+a^2}
\end{align}
であり，等号は$(\cos\theta,\sin\theta)=\dfrac{(t,a)}{\sqrt{t^2+a^2}}$のときに成り立つ．したがって\cref{eq:2}から$|PX|^2$の最小値は
\begin{align}
  |PX|^2\ge(2t^2+a^2+1)-2\sqrt{t^2+a^2}\equiv f(p) \quad (p=t^2\ge0) \label{eq:3}
\end{align}
で与えられる．
以下$a,t$を動かした時の$f(p)$の最小値を考える．
そこで$f(p)$の一階微分を考えると
\begin{align}
  f'(p)=2-\dfrac{1}{\sqrt{p+a^2}}=\dfrac{2\sqrt{p+a^2}-1}{\sqrt{p+a^2}} \label{eq:4}
\end{align}
だから，$a$の値によって場合分けする．

\subparagraph{(i) $a\ge\dfrac{1}{2}$のとき}

任意の$p\ge0$で$\sqrt{p+a^2}\ge a\ge\dfrac{1}{2}$だから，\cref{eq:4}より$f'(p)\ge0$である．よって$f$は$p\ge0$で単調増加で，$p=0\, (t=0)$のとき最小となり
\begin{align}
  \min f(p)=f(0)=a^2-2a+1=(a-1)^2
\end{align}
である．よって\cref{eq:3}より，$|PX|\ge0$に注意して
\begin{align}
  \min|PX|=|a-1|
\end{align}
である．

\subparagraph{(ii) $0<a\le\dfrac{1}{2}$のとき}

\cref{eq:4}より$f'(p)=0\iff p=\dfrac{1}{4}-a^2\,(\ge0)$であり，$f(p)$の増減表は以下のようになる．

\begin{table}[htb]
  \centering
  \caption{$f(p)$ の増減表}
  \begin{tabular}{c|ccccc}
    $p$ & $0$ & $\cdots$ & $\dfrac{1}{4}-a^2$ & $\cdots$ \\
    \hline
    $f'(p)$ & & $-$ & $0$ & $+$ \\
    \hline
    $f(p)$ & & $\searrow$ & 極小 & $\nearrow$
  \end{tabular}
\end{table}

したがって$p=\dfrac{1}{4}-a^2$で$f(p)$は最小となり，
\begin{align}
  \min f(p)
  &=f\left(\dfrac{1}{4}-a^2\right) \\
  &=2\left(\dfrac{1}{4}-a^2\right)+a^2+1-2\sqrt{\dfrac{1}{4}-a^2+a^2} \\
  &=\dfrac{1}{2}-2a^2+a^2+1-1 \\
  &=\dfrac{1}{2}-a^2
\end{align}
である．よって
\begin{align}
  \min|PX|=\sqrt{\dfrac{1}{2}-a^2}
\end{align}
である．
\casesend

以上(i),(ii)より，求める円$C$と直線$l$の最短距離は
\begin{align}
  \min|PX|=
  \begin{cases}
    \sqrt{\dfrac{1}{2}-a^2} & \left(0<a\le\dfrac{1}{2}\right) \\
    |a-1| & \left(\dfrac{1}{2}\le a\right)
  \end{cases}
\end{align}
である．$\cdots$(答)

{\bf[別解]}

[解]と同じ設定のもとで，\cref{eq:1,eq:2}を用いる．
[解]では$t,a$を固定して先に$\theta$を動かしたが，今度は$a,\theta$を固定して先に$t$を動かす．

\cref{eq:2}を$t$について平方完成すると
\begin{align}
  |PX|^2
  & =2\left(t-\dfrac{c}{2}\right)^2-\dfrac{c^2}{2}+a^2-2as+1 \\
  & \ge -\dfrac{c^2}{2}+a^2-2as+1 \quad \left(t=\dfrac{c}{2}\right)\\
  & =\dfrac{s^2}{2}-2as+a^2+\dfrac{1}{2} \\
  & =\dfrac{1}{2}(s-2a)^2-a^2+\dfrac{1}{2} \label{eq:5}
\end{align}
である．$s=\sin\theta$は$-1\le s\le1$を動くから，\cref{eq:5}を最小にするのは，区間$[-1,1]$のうち$2a$に最も近い$s$をとるときである．

\subparagraph{(i) $0<a\le\dfrac{1}{2}$のとき（$2a\le1$）}

$s=2a$が区間$[-1,1]$内にあるからこれを選べて，
\begin{align}
  \min|PX|^2=-a^2+\dfrac{1}{2} \quad \therefore \min|PX|=\sqrt{\dfrac{1}{2}-a^2}
\end{align}
である．

\subparagraph{(ii) $\dfrac{1}{2}\le a$のとき（$2a\ge1$）}

区間$[-1,1]$のうち$2a$に最も近い点は$s=1$だから，\cref{eq:5}に代入して
\begin{align}
  \min|PX|^2=\dfrac{1}{2}(1-2a)^2-a^2+\dfrac{1}{2}=a^2-2a+1=(a-1)^2 \quad \therefore \min|PX|=|a-1|
\end{align}
である．

以上(i),(ii)はいずれも[解]の結果と一致する．$\cdots$(答)

\end{document}
